\begin{samepage}
\subsubsection{Commuting spectral flows and higher-order interactions}
\label{sec:spectral-flows}

The spectral invariants provide Hamiltonians as well as conserved
quantities. Choosing another invariant changes the evolution of the
same mode matrices while preserving their spectral curve. This is
the commuting-flow construction of the AKS framework
\cite{AdlerVanMoerbeke2018}, and applies to arbitrary positive
$L,M,N$.\par
\end{samepage}

For each integer $r\geq2$, define the trace Hamiltonian
\begin{equation}\label{eq:hierarchy-traces}
 H_r=\frac1r\tr\mathsf X^r.
\end{equation}
It is a spectral invariant at $\nu=0$ and is a polynomial of degree
at most $r$ in the modes. By
\cref{prop:canonical-commutation,eq:canonical-trace-H}, these
Hamiltonians commute with one another, with $H$ and $H_0$, and with
all the power charges and spectral coefficients. The complete
integral set in \cref{thm:canonical-independence} is therefore
preserved by every such flow.

For the $H_r$ evolution, the bracket
\cref{eq:canonical-frozen} gives an explicit Lax representation:
\begin{equation}\label{eq:hierarchy-lax}
 i\dot{\mathsf X}=[\mathsf D,\mathsf X^{r-1}],\qquad
 i\dot{\mathsf L}=[\mathsf L,\mathsf T_r],\qquad
 \mathsf T_r(\nu)=
       \frac{\mathsf X^{r-1}-\mathsf L(\nu)^{r-1}}{\nu}.
\end{equation}
The numerator vanishes at $\nu=0$, so $\mathsf T_r$ is polynomial
in $\nu$. Indeed, $[\mathsf L,\mathsf L^{r-1}]=0$ and
$[\mathsf X,\mathsf X^{r-1}]=0$ reduce its commutator to
$[\mathsf D,\mathsf X^{r-1}]$. The block-diagonal entries of
$\mathsf X$ remain fixed because $\mathsf D$ is scalar on each
of its three blocks.

The fourth member $H_4=\tr\mathsf X^4/4$ includes quartic
interactions, with lower-degree terms supplied by the fixed diagonal
entries. Appendix~\ref{sec:quartic-expansion} gives its expanded
quartic part.

For a constant $\alpha$, the combined Hamiltonian and time matrix are
\begin{equation}\label{eq:hierarchy-combination}
 H_\alpha=H+\alpha H_4,\qquad
 \mathsf T_\alpha=\mathsf T+\alpha\mathsf T_4.
\end{equation}
The same initial state and phase coefficients give the same
conserved polynomial $F(\lambda,\nu)$ under either evolution.
In a Jacobian or Prym linearization, if the geometric generators $i\partial_t$ have translation
vectors $\boldsymbol U_H,\boldsymbol U_4$, the divisor's Abel
coordinate evolves as
\begin{equation}\label{eq:hierarchy-direction}
 z(t)=z(0)-it(\boldsymbol U_H+\alpha\boldsymbol U_4).
\end{equation}
This is the spectral interpretation of choosing a different scalar
evolution within the commuting family
\cite{Beauville1990,Vanhaecke2001}.

The hyperelliptic realization in \cref{sec:BA-realization} selects
the scalar Schr\"odinger direction $z(t)=z_0-it e_g$ for $H_0$.
More generally, \cref{thm:BA-higher-flows} identifies the spectral
flows with multipoint BA times and constant reciprocal mode phases
for arbitrary group sizes. On each birational spectral level the
full family supplies every Abel direction.
